\chapter{How to Lower High Blood Pressure}
\index{How to Lower High Blood Pressure}

\section*{Definition and Overview}
Lowering high blood pressure is one of the most effective ways to reduce the risk of heart attack, stroke, kidney disease, and premature death. Blood pressure can be lowered by lifestyle changes, medication, or, most often, a combination of both. Effective strategies include reducing salt, losing weight, increasing physical activity, limiting alcohol, quitting smoking, and managing stress. When lifestyle change is not enough, blood pressure medicines are highly effective and safe. Most people can achieve healthy blood pressure with consistent effort and regular review.

\section*{Epidemiology}
Around one in three adults has high blood pressure, and it is a leading contributor to the cardiovascular disease that causes a quarter of all deaths. Blood pressure control rates have improved with treatment, but many people remain uncontrolled, often because medicines are not taken consistently or doses are not adjusted. Every sustained 10 mmHg reduction in systolic blood pressure lowers the risk of stroke by about a third and heart attack by about a fifth, so even modest improvements matter.

\section*{Aetiology and Pathophysiology}
\begin{itemize}
    \item \textbf{Salt:} excess sodium raises blood volume and pressure; reducing salt lowers pressure within weeks
    \item \textbf{Weight:} excess weight increases the work of the heart; weight loss lowers pressure
    \item \textbf{Physical activity:} regular exercise strengthens the heart and lowers resting pressure
    \item \textbf{Alcohol:} more than two standard drinks daily raises pressure
    \item \textbf{Smoking:} nicotine raises pressure and damages arteries
    \item \textbf{Potassium:} adequate potassium helps balance sodium
    \item \textbf{Stress:} chronic stress activates the nervous system and raises pressure
    \item \textbf{Medication classes:} ACE inhibitors, ARBs, calcium channel blockers, and diuretics lower pressure through different mechanisms
\end{itemize}

\section*{Clinical Features}
\begin{itemize}
    \item Blood pressure targets: usually below 140/90, and often lower in high-risk groups
    \item Lifestyle changes can lower pressure by 5--15 mmHg
    \item Combination therapy: most people need two or more medicines
    \item Home monitoring: self-measured readings guide treatment
    \item Effects of change: blood pressure falls within weeks of lifestyle change
    \item Medication is usually lifelong once started
    \item Review and titration: doses are adjusted until targets are met
\end{itemize}

\section*{Differential Diagnosis}
\begin{itemize}
    \item The approach depends on the level of blood pressure and overall cardiovascular risk
    \item Secondary causes such as kidney disease and hormone disorders are considered in resistant cases
    \item White-coat and masked hypertension are identified by home and ambulatory monitoring
    \item Lifestyle change is appropriate for all; medication is added according to risk
    \item Individual targets vary with age and other conditions
\end{itemize}

\section*{When to Seek Medical Care}
Adults should have blood pressure checked regularly, and anyone with elevated readings should be reviewed and followed up.

\textbf{Contact your local emergency services for:}
\begin{itemize}
    \item Very high readings (above 180/120) with headache, blurred vision, or confusion
    \item Chest pain, breathlessness, or signs of stroke
    \item Sudden weakness, facial drooping, or difficulty speaking
    \item Fainting or collapse
\end{itemize}

\section*{Diagnosis and Investigations}
\begin{itemize}
    \item \textbf{Blood pressure measurement:} repeated readings with correct technique
    \item \textbf{Home and ambulatory monitoring:} confirm the diagnosis and response
    \item \textbf{Cardiovascular risk assessment:} cholesterol, blood sugar, and risk factors
    \item \textbf{Kidney tests:} blood and urine tests
    \item \textbf{ECG:} check for heart strain
    \item \textbf{Regular review:} track progress and adjust treatment
\end{itemize}

\section*{Management}
\begin{itemize}
    \item \textbf{Reduce salt:} aim for under 5--6 g of salt daily; avoid processed foods
    \item \textbf{Lose weight:} even 5--10\% weight loss lowers pressure significantly
    \item \textbf{Exercise:} at least 150 minutes of moderate activity weekly
    \item \textbf{Limit alcohol:} no more than ten standard drinks weekly
    \item \textbf{Quit smoking:} reduces pressure and overall risk
    \item \textbf{Eat well:} the DASH-style diet rich in vegetables, fruit, and wholegrains
    \item \textbf{Medication:} start when lifestyle change is insufficient or risk is high
    \item \textbf{Adherence:} take medicines consistently and attend reviews
\end{itemize}

\section*{Complications}
\begin{itemize}
    \item Untreated hypertension: heart attack, stroke, kidney failure, and vision loss
    \item Medication side effects, usually manageable
    \item Stopping medicines suddenly, which can cause rebound rises
    \item Missed doses, the most common reason for poor control
    \item Interaction of medicines with other drugs
\end{itemize}

\section*{Prognosis and Outcomes}
Lowering blood pressure dramatically reduces cardiovascular risk, and the benefits appear quickly and grow over time. Most people can achieve control with lifestyle change and one to three medicines, and controlled blood pressure restores near-normal risk. Adherence is the strongest predictor of good outcomes, and regular review keeps treatment on track. With consistent management, high blood pressure becomes a well-controlled condition.

\section*{Prevention and Self-Care}
\begin{itemize}
    \item Know your blood pressure numbers
    \item Reduce salt and eat a plant-rich diet
    \item Maintain a healthy weight and waist
    \item Be physically active most days
    \item Limit alcohol and quit smoking
    \item Take medicines exactly as prescribed
    \item Measure blood pressure at home as advised
    \item Attend regular reviews and dose adjustments
\end{itemize}

\section*{Sources and Further Reading}
The following reputable sources provide further information for healthcare students and patients seeking reliable, current advice about lowering blood pressure.

\begin{itemize}
    \item Heart Foundation. \textit{Lowering blood pressure.} https://www.heartfoundation.org.au/
    \item Healthdirect. \textit{High blood pressure: what you can do.} https://www.healthdirect.gov.au/high-blood-pressure
    \item Kidney Health Australia. \textit{Blood pressure and kidneys.} https://www.kidney.org.au/
    \item National Heart Foundation of Australia. \textit{Hypertension management guidelines.} https://www.heartfoundation.org.au/
    \item NHS. \textit{High blood pressure treatment.} https://www.nhs.uk/conditions/high-blood-pressure/
    \item World Health Organization. \textit{Hypertension.} https://www.who.int/
\end{itemize}
